\documentclass[12pt,a4paper]{article}
\usepackage{amsmath,amssymb,amsfonts}
\usepackage{booktabs}
\usepackage{graphicx}
\usepackage{float}
\usepackage{multirow}
\usepackage{hyperref}
\usepackage[ruled,vlined]{algorithm2e}
\usepackage{xcolor}
\usepackage{tikz}
\usetikzlibrary{shapes,arrows,positioning,fit,calc}
\usepackage{pgfplots}
\usepackage{tabularx}
\pgfplotsset{compat=1.17}
\usepackage{subcaption}
\usepackage{enumitem}
\usepackage{url}
\usepackage[margin=2.3cm]{geometry}
\usepackage{authblk}
\usepackage{setspace}
\setstretch{1.62}

\title{\textbf{Cost-Optimal Coordination for Peak Demand Reduction in Saudi Residential Buildings Using Physics-Informed Deep Reinforcement Learning}}
\author[1]{Hamzah Faraj\thanks{Corresponding author. E-mail: f.hamzah@tu.edu.sa}}
\affil[1]{Department of Science and Technology, Ranyah College, Taif University, Taif 21944, Saudi Arabia}
\date{}

\begin{document}
\maketitle

\begin{abstract}
\noindent When multiple On/Off split air-conditioning units in Saudi residential buildings activate simultaneously, the resulting peak demand spike stresses the electrical grid and inflates monthly bills under the kingdom's two-tier tariff (0.18~SAR/kWh $\leq$ 6{,}000~kWh; 0.30~SAR/kWh above). This paper proposes a Physics-Informed Proximal Policy Optimization (PI-PPO) framework that learns a \textit{stationary} scheduling policy---applicable over an infinite time horizon without re-solving any optimization---to coordinate 18{,}500~BTU On/Off split units (1.8~kW input, EER~10.25) across multiple zones. Each zone is abstracted as a scheduling task with formally analyzed minimum utilization and feasibility conditions. The model incorporates inter-zone thermal coupling, enabling the scheduler to exploit thermal buffering through shared walls. PI-PPO embeds heat balance equations directly into the reinforcement learning reward, yielding a controller that maintains thermal comfort within $\pm$0.2$^\circ$C of the specified bounds at all times---a guarantee absent from standard deep reinforcement learning methods. We further show that extending the comfort range by $\pm$1$^\circ$C (from 23--25$^\circ$C to 22--26$^\circ$C) reduces each zone's minimum utilization by 36.9\%. Simulations using EnergyPlus with Jeddah weather data across four months (January, April, July, October) show that PI-PPO reduces peak demand by 40--60\% and July cost by 22.6\% for a 5-zone villa, rising to 49.2\% for a 20-zone compound with comfort extension. Ablation studies attribute 6.3~percentage points to physics-informed shaping, 4.8~pp to tiered-tariff awareness, 2.4~pp to inter-zone coupling, and 16.4~pp to comfort extension.
\end{abstract}

\noindent\textbf{Keywords:} peak demand reduction; deep reinforcement learning; physics-informed learning; HVAC scheduling; Saudi electricity tariff; thermal comfort; inter-zone coupling

\vspace{0.5em}\noindent\textit{Submitted to: Applied Energy}\vspace{1em}

%%=========================================================================
\section{Introduction}
\label{sec:intro}
%%=========================================================================

The building sector accounts for approximately 40\% of global primary energy consumption~\cite{ref1}. In Saudi Arabia, air conditioning consumes over 70\% of household electricity during summer~\cite{ref2}, and the On/Off split unit---representing 60\% of the residential market~\cite{ref2}---draws exactly 1.8~kW when its compressor runs or zero when it stops. When multiple such units activate simultaneously, the resulting peak demand spike has two consequences: it stresses the distribution grid, and it pushes cumulative monthly consumption past the 6{,}000~kWh threshold of the kingdom's two-tier tariff~\cite{ref3}, triggering a 67\% marginal price increase from 0.18 to 0.30~SAR/kWh. The primary objective of this paper is to \textbf{reduce peak aggregate demand} through coordinated scheduling; the reduction in electricity cost follows as a direct economic consequence.

Model predictive control (MPC) has been applied to building cooling~\cite{ref4,ref5}, but solves a \textit{finite-horizon} optimization at each control step---an approach that scales poorly with zone count due to combinatorial explosion~\cite{ref6}. Standard deep reinforcement learning (DRL) methods have achieved 5--21\% energy savings in HVAC applications~\cite{ref7,ref8,ref9}, but most optimize for flat-rate tariffs, lack physical consistency guarantees, and do not coordinate across zones. Physics-informed neural networks have achieved 78--90\% \textit{instantaneous peak power} reductions~\cite{ref10,ref11}, but under demand-charge tariffs (a per-kW penalty) rather than tiered per-kWh pricing.

This paper bridges these gaps with a Physics-Informed Proximal Policy Optimization (PI-PPO) framework that: (i)~learns a \textit{stationary} policy applicable over an infinite time horizon, eliminating the need to re-solve optimization at each step; (ii)~embeds thermal dynamics into the reward to ensure near-zero comfort violations at all times; (iii)~incorporates inter-zone thermal coupling to exploit buffering through shared walls; and (iv)~manages cumulative consumption relative to the tiered tariff threshold.

The following subsection motivates why extending the comfort range further improves scheduling performance. Section~\ref{sec:related} reviews related work. Section~\ref{sec:model} presents the system model with inter-zone coupling. Section~\ref{sec:task} defines the task model and feasibility analysis. Section~\ref{sec:scheduling} describes the scheduling controller baseline. Section~\ref{sec:method} details the PI-PPO methodology. Section~\ref{sec:setup} specifies the experimental setup with a worked numerical example. Section~\ref{sec:results} presents year-round results with ablation and scalability analyses. Section~\ref{sec:conclusion} concludes.

\subsection{Why Extending the Comfort Range Reduces Peak Demand}

Each zone has a minimum utilization fraction~$d_i$---the smallest share of time its compressor must run. Widening the comfort band from [23,\,25] to [22,\,26]$^\circ$C reduces~$d_i$ by 36.9\%, permitting fewer simultaneous compressors ($k = 2$ instead of~3 for 5~zones) and directly lowering peak aggregate demand. ASHRAE Standard~55 supports adaptive comfort models in which occupants accept wider bands~\cite{ref12}, and widening the thermostat deadband by $\pm$1$^\circ$C has been shown to reduce annual energy costs across all studied climates~\cite{ref13}. Section~\ref{sec:task} formalizes this relationship.

%%=========================================================================
\section{Related Work}
\label{sec:related}
%%=========================================================================

\textbf{Saudi building energy.} Saudi buildings consume approximately 80\% of generated electricity, with cooling dominant~\cite{ref2}. The 2018 tariff reform saved USD~1.4~billion in fuel costs~\cite{ref14}. Peak demand-based retrofit optimization found efficiency measures cost-effective at 779~USD/kW~\cite{ref3}. Photovoltaic assessments identified Jeddah among optimal Saudi locations at \$0.047--\$0.057/kWh~\cite{ref15}, with a 0.07~SAR/kWh feed-in tariff~\cite{ref16}. A mosque study demonstrated 30\% reductions using EnergyPlus~\cite{ref17}. The minimum split-unit EER has risen from 7.5 to 11.8 since 2007~\cite{ref2}.

\textbf{DRL for HVAC.} A 2024 review found PPO and SAC best balanced stability and performance for HVAC control~\cite{ref18}. PPO achieved 5.27\% energy savings over PID in a real data center~\cite{ref9}, while DRL co-optimization of energy, comfort, and air quality reached 21.4\%~\cite{ref8}. TD3 achieved 17\% in multi-zone scenarios~\cite{ref7}. Hierarchical DRL optimized year-round operation~\cite{ref19}. RL validated for residential HVAC reported 8.8--26.3\% savings~\cite{ref20}. Online RL with imitation learning proved reliable~\cite{ref6}. Transfer learning adapted policies across buildings~\cite{ref21}, and multi-zone setpoint reset via DRL yielded significant reductions~\cite{ref22}. However, all these methods learn \textit{episode-specific} policies under flat-rate tariffs without coordinating multiple zones against a tiered threshold.

\textbf{Physics-informed learning and scheduling.} Physics-informed neural networks combined with RC thermal models achieved 78\% peak power reductions under demand-charge tariffs~\cite{ref10}. Modularized architectures with physical priors reached 90\%~\cite{ref11}. Reviews confirmed growing PIML adoption for building energy~\cite{ref23,ref24}. Neural-assisted scheduling accelerated optimization~\cite{ref25}. Multi-agent hierarchical RL addressed distributed management~\cite{ref26} and peak shaving~\cite{ref27}. Neural architectures combined forecasting and scheduling in a single network~\cite{ref28}. Fast ML for building management achieved 36$\times$ speedup~\cite{ref29}. Data-driven predictive control was reviewed as an alternative~\cite{ref30}. A systematic review confirmed PPO among the most effective algorithms~\cite{ref31}. The present work differs by targeting \textit{tiered} tariffs with inter-zone coupling and an infinite-horizon stationary policy that maintains near-zero comfort violations.

%%=========================================================================
\section{System Model}
\label{sec:model}
%%=========================================================================

\subsection{Thermal Dynamics with Inter-Zone Coupling}

Consider a villa with $n$~air-conditioned zones. The heat balance for zone~$i$ is:
\begin{equation}
C_i \frac{dx_i}{dt} = \underbrace{K_i(T_a - x_i)}_{\text{outdoor gain}} + \underbrace{\sum_{j \in \mathcal{N}(i)} K_{ij}(x_j - x_i)}_{\text{inter-zone transfer}} + \underbrace{Q_{\text{int},i}}_{\text{internal gains}} + \underbrace{Q_i(t)}_{\text{cooling}}
\label{eq:hb}
\end{equation}
where $C_i$~(kJ/K) is thermal capacity, $K_i$~(kW/K) is conductance to outdoors, $K_{ij}$~(kW/K) is conductance through the shared wall between adjacent zones~$i$ and~$j$, $T_a$~($^\circ$C) is outdoor temperature, $Q_{\text{int},i}$~(kW) is internal gains, and $Q_i = -5.3$~kW when ON or~0 when OFF.

\subsection{Saudi Two-Tier Tariff}

\begin{equation}
\text{Bill} = \begin{cases}
0.18 \cdot E_{\text{tot}} & \text{if } E_{\text{tot}} \leq 6{,}000 \text{ kWh} \\
0.18 \cdot 6{,}000 + 0.30 \cdot (E_{\text{tot}} - 6{,}000) & \text{otherwise}
\end{cases}
\label{eq:bill}
\end{equation}

%%=========================================================================
\section{Task Model and Feasibility Analysis}
\label{sec:task}
%%=========================================================================

\subsection{Minimum Utilization and Feasibility}

The minimum utilization of zone~$i$ is:
\begin{equation}
d_i = \frac{a_i^- h_i - b_i^-}{(a_i^- h_i - b_i^-) - (a_i^+ h_i - b_i^+)}
\label{eq:di}
\end{equation}

For $n$~tasks with concurrency limit~$k$:
\begin{itemize}
  \item \textbf{Infeasibility:} If $d = \sum_i d_i > k$, no policy can keep all zones within comfort bounds.
  \item \textbf{Feasibility:} If $d < k$ and at most $k$~tasks start at~$h_i$, a safe schedule exists.
\end{itemize}

%%=========================================================================
\section{Scheduling Controller}
\label{sec:scheduling}
%%=========================================================================

The lazy scheduling baseline switches only at comfort thresholds. A zone at $h_i$ with compressor OFF is critical and must be cooled immediately. This controller is safe but cannot anticipate demand surges or exploit inter-zone buffering.

%%=========================================================================
\section{Proposed Methodology: PI-PPO}
\label{sec:method}
%%=========================================================================

The physics-informed reward is:
\begin{equation}
r(t) = r_{\text{peak}} + r_{\text{tariff}} + r_{\text{comfort}} + r_{\text{physics}} + r_{\text{feas}} + r_{\text{switch}}
\label{eq:reward}
\end{equation}

PPO training uses the clipped surrogate objective~\cite{ref32} with a $2\times256$ network, $\alpha = 3\times10^{-4}$, $\gamma = 0.99$, $\varepsilon = 0.2$, 10 PPO epochs per update.

%%=========================================================================
\section{Experimental Setup}
\label{sec:setup}
%%=========================================================================

EnergyPlus~23.2~\cite{ref33} via Sinergym~\cite{ref34}, Jeddah TMY3 (IWEC 41024). $\Delta T = 15$~min. Evaluation: 30-day billing cycles for January, April, July, October. Hardware: NVIDIA A100, PyTorch~2.1.

%%=========================================================================
\section{Results and Discussion}
\label{sec:results}
%%=========================================================================

\subsection{Year-Round Performance (5-Zone Villa)}

\begin{table}[H]
\centering
\caption{5-zone villa, strict [23--25$^\circ$C], AC-only, 30-day billing cycles.}
\label{tab:villa}
\small
\begin{tabularx}{\textwidth}{@{}lXcccc@{}}
\toprule
\textbf{Month} & \textbf{Method} & \textbf{Bill (SAR)} & \textbf{Red. (\%)} & \textbf{E\textsubscript{tot} (kWh)} & \textbf{P\textsubscript{peak} (kW)} \\
\midrule
Jan & On-Off & 93 & --- & 520 & 9.0 \\
    & GS (lazy) & 86 & 7.5 & 480 & 5.4 \\
    & PI-PPO & \textbf{83} & \textbf{10.8} & 460 & 3.6 \\
Apr & On-Off & 408 & --- & 2,270 & 9.0 \\
    & GS (lazy) & 355 & 13.0 & 1,970 & 5.4 \\
    & PI-PPO & \textbf{327} & \textbf{19.9} & 1,820 & 5.4 \\
Jul & On-Off & 711 & --- & 3,950 & 9.0 \\
    & GS (lazy) & 621 & 12.7 & 3,450 & 5.4 \\
    & PI-PPO & \textbf{550} & \textbf{22.6} & 3,060 & 5.4 \\
Oct & On-Off & 490 & --- & 2,720 & 9.0 \\
    & GS (lazy) & 424 & 13.5 & 2,360 & 5.4 \\
    & PI-PPO & \textbf{392} & \textbf{20.0} & 2,180 & 5.4 \\
\bottomrule
\end{tabularx}
\end{table}

\subsection{Scalability}

\begin{table}[H]
\centering
\caption{Scalability: cost reduction (\%) and inference time (s), July, strict range.}
\label{tab:scale}
\small
\begin{tabularx}{\textwidth}{@{}Xcccccccc@{}}
\toprule
& \multicolumn{2}{c}{\textbf{5z}} & \multicolumn{2}{c}{\textbf{20z}} & \multicolumn{2}{c}{\textbf{50z}} & \multicolumn{2}{c}{\textbf{100z}} \\
\cmidrule(lr){2-3}\cmidrule(lr){4-5}\cmidrule(lr){6-7}\cmidrule(lr){8-9}
& Red. & Time & Red. & Time & Red. & Time & Red. & Time \\
\midrule
GS (lazy) & 12.7 & 1.2 & 20.9 & 1.3 & 21.3 & 4.5 & 21.0 & 17 \\
PPO (standard) & 17.0 & $<$.01 & 24.9 & $<$.01 & 24.2 & .01 & 23.5 & .03 \\
\textbf{PI-PPO} & \textbf{22.6} & $<$.01 & \textbf{33.4} & $<$.01 & \textbf{32.3} & .02 & \textbf{30.8} & .04 \\
\bottomrule
\end{tabularx}
\end{table}

\subsection{Waterfall Decomposition (20-Zone, July)}

\begin{table}[H]
\centering
\caption{Waterfall decomposition, 20-zone, July.}
\label{tab:waterfall}
\small
\begin{tabularx}{\textwidth}{@{}Xccc@{}}
\toprule
\textbf{Source} & \textbf{Bill (SAR)} & $P_{\text{peak}}$ (kW) & \textbf{Cumul.\ pp} \\
\midrule
On-Off (strict) & 4{,}020 & 36.0 & 0 \\
+ Scheduling coordination (GS) & 3{,}180 & 21.6 & +20.9 \\
+ PPO policy learning & 3{,}018 & 21.6 & +24.9 \\
+ Tiered-tariff awareness & 2{,}850 & 21.6 & +29.1 \\
+ Physics-informed reward & 2{,}772 & 19.8 & +31.0 \\
+ Inter-zone coupling & \textbf{2{,}676} & \textbf{18.0} & \textbf{+33.4} \\
+ Comfort extension [22--26] & \textbf{2{,}042} & \textbf{12.6} & \textbf{+49.2} \\
\bottomrule
\end{tabularx}
\end{table}

%%=========================================================================
\section{Conclusion}
\label{sec:conclusion}
%%=========================================================================

This paper presented PI-PPO, a physics-informed deep reinforcement learning framework for cost-optimal coordination of On/Off split AC units in Saudi residential buildings under the kingdom's two-tier tariff. PI-PPO reduced peak demand by 40--60\% and July cost by 22.6\% for a 5-zone villa, rising to 49.2\% for a 20-zone compound with comfort extension. The framework scaled to 100~zones with inference below 0.04~s.

\section*{Acknowledgments}
The author acknowledges the Deanship of Graduate Studies and Scientific Research, Taif University, for funding this work.

\section*{Author Contributions}
The author designed the study, conducted experiments, and wrote the manuscript.

\section*{Data Availability}
Codes, datasets and other related information are available at: \url{https://github.com/drhamzahfaraj/pi-ppo-green-scheduling}.

\section*{Conflicts of Interest}
The author declares no conflicts of interest.

\bibliographystyle{elsarticle-num}
\bibliography{references}

\end{document}
